% Based on Template from IIS, ETHz
%%%%%%%%%%%%%%%%%%%%%%%%%%%%%%%%%%%%%%%%%%%%%%%%%%%%%%%%%%%%%%%%%%%%%%%
%%%%%%%%%%%%%%%%%%%%%%%%%%%%%%%%%%%%%%%%%%%%%%%%%%%%%%%%%%%%%%%%%%%%%%%
%%%%%                                                                 %
%%%%%     <file_name>.tex                                             %
%%%%%                                                                 %
%%%%% Author:      <author>                                           %
%%%%% Created:     <date>                                             %
%%%%% Description: <description>                                      %
%%%%%                                                                 %
%%%%%%%%%%%%%%%%%%%%%%%%%%%%%%%%%%%%%%%%%%%%%%%%%%%%%%%%%%%%%%%%%%%%%%%
%%%%%%%%%%%%%%%%%%%%%%%%%%%%%%%%%%%%%%%%%%%%%%%%%%%%%%%%%%%%%%%%%%%%%%%

\chapter{Introduction}
Give an overview of the problem, and put your work into a bigger
context. Motivate the questions addressed in this work and summarize
your contributions. 

If you want to \emph{highlight} a word or term use \verb|\emph{}|.

If you work with many acronyms we recommend the use of an acronym library. You can define acronyms in the lower part of the glossaries.tex file. When first using the acronym it will be introduced \gls{pbl}, after that the shortform will be displayed \gls{pbl}.

\gls{scuba} diving is a recreational activity that involves breathing one or more compressed gas mixes.
Dive computers in their most basic form consist of a depth sensor, a real-time clock, processing unit, and a display. 
Their primary function is to display dive time and depth, as well as to calculate and display decompression obligation and \gls{otu}s during the dive. 
Furthermore, most computers show statistics about recent dives or allow downloading the dive logs for further analysis and accurate logging when on the surface.

\section{Motivation}
Dive Computers commonly sample depth and update the corresponding display multiple times per second, while log rate (the stored logs for later retrieval) is usually user-defined between 1 and 30 seconds. Recomputation of decompression schedules cannot be as simply observed; thus, it can only assumed that it happens on a scale of at least every few seconds. 

As Prof Dr Simon Mitchell explains, hyperbaric medicine and especially decompression science is 'bucket chemistry', meaning that the extent to which accurate O2 time and decompression tracking is useful is on the order of seconds. However, this does not mean that more frequent sampling is useless for all applications, most divers rely on their dive computers to refresh the depth display in increments of $\SI{10}{\cm}$ with a sub-second delay. 

However, there is no indication that knowledge about typical dive profiles influences these rates, i.e., there is no research in the public domain that investigates the feasibility of a variable refresh rate for all the previously listed components. 

\section{Objective}
what was the objective of your thesis? also add your research questions here

\subsection{Research Questions}
\begin{enumerate}
\item Is it possible to create a fully Open-Source dive computer which can use both exponential and linear-exponential algorithms, based on the user's preference?
\item Can an appropriately-sized \gls{pcb} be designed and still be hand-soldered?
\item Can sufficiently much energy and log space be saved by employing task variable-rate task scheduling / procedure running to offset the incurred overhead by the variable-rate algorithms?
\end{enumerate}

\subsection{Outline}
First, an introduction to measurement-based embedded devices, energy-saving measures in battery-power devices, and basic hyperbaric medicine theory is presented, to the extent that it is relevant to divers and dive computer developers. 
Next, we look at related work to save energy in embedded devices, mainly by running at tasks variable rate. 
Following this, the specific implementation of this project in hardware and software is described. 
Chapter 5 will collect and compare the results received using different sortware.
Finally, the results will be critically discussed and compared with those in the related work.
